\documentclass[aspectratio=169]{beamer}

% ======================
% THEME & LANG
% ======================
\usetheme{Madrid}
\usecolortheme{default}
\useinnertheme{rounded}
\usefonttheme{professionalfonts}

\usepackage[spanish]{babel}
\usepackage[utf8]{inputenc}
\usepackage[T1]{fontenc}
\usepackage{lmodern}

% ======================
% PACKAGES
% ======================
\usepackage{graphicx}
\usepackage{booktabs}
\usepackage{tabularx}
\usepackage{array}
\usepackage{multicol}
\usepackage{pgfplots}
\pgfplotsset{compat=1.18}
\usepackage{pgf-pie}
\usepackage{tikz}
\usetikzlibrary{arrows.meta,positioning,shapes.geometric,fit,calc}

% ======================
% COLORES PERSONALIZADOS
% ======================
% UNAM branding colors
\definecolor{UNAMBlue}{RGB}{0,38,84}
\definecolor{UNAMGold}{RGB}{255,204,0}
\definecolor{ForestGreen}{RGB}{34,139,34}
\definecolor{Orange}{RGB}{255,140,0}
\definecolor{UNAMRed}{RGB}{171,39,39}

% ======================
% META
% ======================
\title[Avances Tesis - Nov 2025]{Modelado Evolutivo del Crecimiento Urbano}
\subtitle{Integración WoE, Autómatas Celulares y Algoritmos Genéticos \\ \textbf{Avances Noviembre 2025}}
\author{Rodrigo Moreno López}
\institute{Maestría en Ciencias de la Computación \\ Universidad Nacional Autónoma de México}
\date{6 de Diciembre de 2025 \\ Coloquio de Investigación II - Tercer Semestre}

% ======================
% MACROS
% ======================
\newcommand{\TotalRefs}{67}
\newcommand{\RefsCA}{15}
\newcommand{\RefsCalib}{12}
\newcommand{\RefsRS}{18}
\newcommand{\RefsImpacto}{8}
\newcommand{\RefsTools}{6}
\newcommand{\RefsWoE}{8}

% Safe derived values for PGFPlots
\pgfmathtruncatemacro{\RefsOthers}{\TotalRefs-\RefsCalib}
\pgfmathtruncatemacro{\YMax}{\TotalRefs+5} 

% ======================
% DOCUMENTO
% ======================
\begin{document}

% TÍTULO
\begin{frame}
  \titlepage
\end{frame}

% AGENDA
\begin{frame}{Agenda}
  \tableofcontents[hideallsubsections]
\end{frame}

\section{Contexto del Proyecto}
\begin{frame}{Contexto y Problema de Investigación}
  \begin{itemize}
    \item El crecimiento urbano en México ejerce presión sobre ecosistemas, infraestructura y calidad de vida.
    \item Los \textbf{Autómatas Celulares (AC)} son efectivos para capturar dinámicas espaciales a partir de reglas locales.
    \item Reto: \textbf{calibración} de parámetros (pesos de influencia, umbrales, factores espaciales) de forma reproducible y no sesgada.
    \item \textcolor{UNAMBlue}{\textbf{Propuesta innovadora:}} \textbf{WoE + AC + AG} con datos satelitales históricos de Querétaro (1984-2020).
  \end{itemize}
  
  \vspace{0.3cm}
  \begin{alertblock}{Caso de Estudio: Querétaro}
  \small 37 imágenes históricas de Google Earth • Crecimiento urbano 340\% (1984-2020) • 
  Metodología transferible a otras ciudades mexicanas
  \end{alertblock}
\end{frame}

\section{Objetivos y Alcance}
\begin{frame}{Objetivo General}
  \Large
  \textcolor{UNAMBlue}{Desarrollar un \textbf{modelo evolutivo WoE-AC-AG} para simular y predecir el crecimiento urbano de Querétaro, optimizando parámetros con datos satelitales históricos y validación temporal robusta.}
\end{frame}

\begin{frame}{Objetivos Específicos}
  \begin{enumerate}
    \item \textbf{Pipeline de datos:} Procesar y clasificar 37 imágenes satelitales históricas (1984-2020)
    \item \textbf{Weight of Evidence:} Implementar sistema WoE para cuantificar relaciones espaciales
    \item \textbf{Autómata Celular:} Desarrollar AC con reglas de transición basadas en WoE
    \item \textbf{Optimización evolutiva:} Integrar AG para calibración automática de parámetros
    \item \textbf{Validación temporal:} Evaluar con métricas espaciales (IoU, FOM, Kappa) en múltiples períodos
    \item \textbf{Escenarios futuros:} Explorar proyecciones de crecimiento 2025-2050
  \end{enumerate}
\end{frame}

\section{Avances Metodológicos}

\begin{frame}{Arquitectura del Dataset Desarrollado}
  \begin{columns}
    \begin{column}{0.6\textwidth}
      \textbf{Resumen:}
      \begin{itemize}
        \item 37 imágenes históricas (1984-2020)
        \item Clasificación binaria urbano/no-urbano
        \item Área de estudio fija, resolución constante
      \end{itemize}
    \end{column}
    \begin{column}{0.4\textwidth}
      \begin{center}
        \begin{tikzpicture}[scale=0.9]
          % Línea temporal simple
          \draw[UNAMBlue, very thick] (0,0) -- (7,0);
          % marcas principales
          \foreach \x/\year in {0/1984,2/2000,4/2010,6/2020} {
            \draw[UNAMBlue, thick] (\x,-0.15) -- (\x,0.15);
            \node[below, font=\scriptsize, text=UNAMBlue] at (\x,-0.2) {\year};
          }
          % título años
          \node[above right, font=\footnotesize, text=UNAMRed] at (0.05,0.2) {37 años};
        \end{tikzpicture}
      \end{center}
    \end{column}
  \end{columns}
\end{frame}

\begin{frame}{Pipeline de Procesamiento Implementado}
  \begin{center}
  \begin{tikzpicture}[
      node distance=5mm and 5mm,
      every node/.style={font=\scriptsize},
      box/.style={
        draw, rounded corners=2mm, align=center, inner sep=2mm, fill=gray!10,
        text width=2.2cm, minimum height=1.2cm
      },
      greenbox/.style={box, fill=ForestGreen!20},
      bluebox/.style={box, fill=UNAMBlue!20, text=white},
      >={Latex},
      scale=0.8, transform shape
    ]
    
    % FILA 1: Adquisición y Preprocesamiento
    \node[greenbox] (raw) {\textbf{Imágenes Crudas}\\Google Earth\\37 años históricos\\RGB 8-bit};
    \node[greenbox, right=of raw] (geo) {\textbf{Georreferenciación}\\Coordenadas exactas\\Superposición píxel\\Control calidad};
    \node[greenbox, right=of geo] (norm) {\textbf{Normalización}\\Temporal consistente\\Período seco nov-abr\\Sin nubes/sombras};
    
    % FILA 2: Extracción de Características  
    \node[bluebox, below=of raw] (lbp) {\textbf{Análisis LBP}\\Texturas urbanas\\P=8, R=1.0\\Patrones locales};
    \node[bluebox, right=of lbp] (sobel) {\textbf{Filtros Sobel}\\Detección bordes\\Estructuras lineales\\Carreteras};
    \node[bluebox, right=of sobel] (cluster) {\textbf{K-means}\\2-4 clusters\\Clasificación\\no supervisada};
    
    % FILA 3: Clasificación y Validación
    \node[box, fill=Orange!20, below=of lbp] (svm) {\textbf{Clasificación SVM}\\Kernel lineal/RBF\\Entrenamiento\\supervisado};
    \node[box, fill=Orange!20, right=of svm] (binary) {\textbf{Mapas Binarios}\\Urbano/No-urbano\\Validación temporal\\Métricas espaciales};
    \node[box, fill=Orange!20, right=of binary] (woe) {\textbf{Cálculo WoE}\\Variables explicativas\\Information Value\\Discretización};

    % CONEXIONES
    \draw[->, thick] (raw) -- (geo);
    \draw[->, thick] (geo) -- (norm);
    \draw[->, thick] (raw) -- (lbp);
    \draw[->, thick] (geo) -- (sobel);
    \draw[->, thick] (norm) -- (cluster);
    \draw[->, thick] (lbp) -- (svm);
    \draw[->, thick] (sobel) -- (binary);
    \draw[->, thick] (cluster) -- (woe);
  \end{tikzpicture}
  \end{center}
  
  \begin{alertblock}{Estado de Implementación}
  \textcolor{ForestGreen}{\textbf{$\checkmark$ Pipeline completo:}} 37 imágenes procesadas • 
  \textcolor{ForestGreen}{\textbf{$\checkmark$ Precisión >99\%}} clasificación • 
  \textcolor{UNAMBlue}{\textbf{$\checkmark$ Sistema WoE}} funcional
  \end{alertblock}
\end{frame}

\begin{frame}{Resultados: Evolución Urbana Cuantificada}
  \begin{columns}
    \begin{column}{0.5\textwidth}
      \textbf{Métricas de Crecimiento:}
      \begin{itemize}
        \item \textbf{1984:} 2,847 hectáreas (línea base)
        \item \textbf{1990:} 3,421 ha (\textcolor{ForestGreen}{+20.2\%} en 6 años)
        \item \textbf{2000:} 5,689 ha (\textcolor{Orange}{+66.3\%} en 10 años)
        \item \textbf{2010:} 8,234 ha (\textcolor{UNAMBlue}{+44.7\%} en 10 años)  
  \item \textbf{2020:} 12,523 ha (\textcolor{UNAMRed}{+52.1\%} en 10 años)
      \end{itemize}
      
      \vspace{0.3cm}
      \begin{exampleblock}{Hallazgo Clave}
      \small Crecimiento \textbf{340\%} total (36 años) • 
      Tasa máxima \textbf{1990-2000} • 
      Sostenimiento \textbf{alto} post-2000
      \end{exampleblock}
    \end{column}
    \begin{column}{0.5\textwidth}
      \begin{center}
      \begin{tikzpicture}[scale=0.8]
        \begin{axis}[
          width=5cm, height=3cm,
          xlabel={\small Año},
          ylabel={\small Hectáreas},
          xmin=1980, xmax=2025,
          ymin=0, ymax=14000,
          grid=major,
          legend pos=north west
        ]
        \addplot[UNAMBlue, thick, mark=*] coordinates {
          (1984, 2847)
          (1990, 3421) 
          (2000, 5689)
          (2010, 8234)
          (2020, 12523)
        };
        \legend{Área urbana}
        \end{axis}
      \end{tikzpicture}
      \end{center}
      \vspace{0.2cm}
      {\small \textbf{Crecimiento 340\%} en 36 años}
    \end{column}
  \end{columns}
\end{frame}

\section{Implementación WoE-AC-AG}

\begin{frame}{Sistema Weight of Evidence Desarrollado}
  \begin{columns}
    \begin{column}{0.6\textwidth}
      \textbf{Fundamento Teórico:}
      $$WoE_i = \ln\left(\frac{N_{1i}/N_1}{N_{0i}/N_0}\right)$$
      
      \textbf{Information Value:}
      $$IV = \sum WoE_i \times \left(\frac{N_{1i}}{N_1} - \frac{N_{0i}}{N_0}\right)$$
      
      \textbf{Capacidades Implementadas:}
      \begin{itemize}
        \item \textcolor{ForestGreen}{\textbf{$\checkmark$}} Discretización automática variables continuas
        \item \textcolor{ForestGreen}{\textbf{$\checkmark$}} Soporte variables categóricas  
        \item \textcolor{ForestGreen}{\textbf{$\checkmark$}} Cálculo Information Value predictivo
        \item \textcolor{ForestGreen}{\textbf{$\checkmark$}} Framework probabilístico bayesiano
      \end{itemize}
    \end{column}
    \begin{column}{0.4\textwidth}
      \textbf{Módulo Técnico:}
      \begin{itemize}
        \item Clase \texttt{WoECalculator}
        \item Integración con pipeline satelital
        \item Variables explicativas: topografía, carreteras, centros urbanos
        \item Base matemática sólida para AC
      \end{itemize}
      
      \vspace{0.3cm}
      \begin{alertblock}{Estado}
      \textcolor{ForestGreen}{\textbf{100\% Implementado}} \\
      Probado con datos Querétaro
      \end{alertblock}
    \end{column}
  \end{columns}
\end{frame}

\begin{frame}{Modelo de Autómata Celular Desarrollado}
  \textbf{Función de Transición Implementada:}
  
  $$P_{i,j}(t+1) = P_{base} \times \prod_{k=1}^{n} e^{W_k \times WoE_{k,i,j}} \times N_{i,j}(t) \times R_{i,j}$$
  
  \vspace{0.3cm}
  \begin{columns}
    \begin{column}{0.5\textwidth}
      \textbf{Componentes:}
      \begin{itemize}
        \item $P_{base}$: Probabilidad base crecimiento
        \item $W_k$: Pesos variables (optimizar con AG)
        \item $WoE_{k,i,j}$: Evidencia espacial celda $(i,j)$  
        \item $N_{i,j}(t)$: Efecto vecindad (Moore 3×3)
        \item $R_{i,j}$: Restricciones (pendiente, áreas protegidas)
      \end{itemize}
    \end{column}
    \begin{column}{0.5\textwidth}
      \textbf{Estado de Desarrollo:}
      \begin{itemize}
        \item \textcolor{ForestGreen}{\textbf{$\checkmark$}} Motor simulación temporal
        \item \textcolor{ForestGreen}{\textbf{$\checkmark$}} Reglas transición probabilísticas
        \item \textcolor{ForestGreen}{\textbf{$\checkmark$}} Integración restricciones espaciales
        \item \textcolor{UNAMBlue}{\textbf{75\%}} Interface con AG (en desarrollo)
      \end{itemize}
      
      \vspace{0.2cm}
      \textbf{Especificaciones:}
      \begin{itemize}
        \item Grid 100m × 100m
        \item Estados: NO\_URBAN, URBAN, RESTRICTED
        \item Vecindad Moore configurable
      \end{itemize}
    \end{column}
  \end{columns}
\end{frame}

\begin{frame}{Sistema de Optimización con Algoritmos Genéticos}
  \begin{columns}
    \begin{column}{0.6\textwidth}
      \textbf{Framework DEAP Implementado:}
      \begin{itemize}
        \item Clase \texttt{WoEACOptimizer} funcional
        \item Optimización multi-objetivo (IoU, Kappa, FoM)
        \item Pipeline integrado \texttt{WoEACPipeline}
        \item 8 parámetros optimizables simultáneamente
      \end{itemize}
      
      \vspace{0.3cm}
      \textbf{Parámetros a Optimizar:}
      \begin{enumerate}
        \item Pesos WoE por variable explicativa (×5)
        \item Factor influencia vecinal ($\alpha$)  
        \item Probabilidad base ($P_{base}$)
        \item Umbral probabilidad transición
      \end{enumerate}
    \end{column}
    \begin{column}{0.4\textwidth}
      \textbf{Configuración Evolutiva:}
      \begin{itemize}
        \item Población: 50 individuos
        \item Generaciones: hasta convergencia (~47)
        \item Cruce: uniforme (prob. 0.7)
        \item Mutación: gaussiana (prob. 0.2)
        \item Selección: torneo (size=3)
      \end{itemize}
      
      \vspace{0.3cm}
      \begin{alertblock}{Estado}
      \textcolor{UNAMBlue}{\textbf{95\% Implementado}} \\
      Calibración en progreso
      \end{alertblock}
    \end{column}
  \end{columns}
\end{frame}

\section{Arquitectura de Software}

\begin{frame}{Sistema Modular Desarrollado}
  \begin{center}
  \begin{tikzpicture}[
      node distance=8mm and 8mm,
      every node/.style={font=\scriptsize},
      module/.style={
        draw, rounded corners=3mm, align=center, inner sep=3mm, 
        text width=3cm, minimum height=1.5cm
      },
      >={Latex},
      scale=0.75, transform shape
    ]
    
    % MÓDULOS PRINCIPALES
    \node[module, fill=ForestGreen!20] (pipeline) {
      \textbf{Pipeline Module}\\
      \texttt{SatelliteImagePipeline}\\
      LBP + Sobel + K-means\\
      Clasificación SVM
    };
    
    \node[module, fill=UNAMBlue!20, text=white, right=of pipeline] (woe) {
      \textbf{WoE Module}\\
      \texttt{WoECalculator}\\
      Discretización automática\\
      Information Value
    };
    
    \node[module, fill=Orange!20, right=of woe] (ca) {
      \textbf{CA Module}\\
      \texttt{WoECellularAutomaton}\\
      Reglas transición\\
      Simulador temporal
    };
    
  \node[module, fill=UNAMRed!20, text=white, below=of woe] (ga) {
      \textbf{GA Module}\\
      \texttt{WoEACOptimizer}\\
      Framework DEAP\\
      Multi-objetivo
    };
    
    \node[module, fill=gray!30, left=of ga] (utils) {
      \textbf{Utils Module}\\
      I/O datos\\
      Visualización\\
      Configuración YAML
    };
    
    \node[module, fill=gray!30, right=of ga] (eval) {
      \textbf{Eval Module}\\
      Métricas espaciales\\
      IoU, Kappa, FoM\\
      Validación temporal
    };

    % CONEXIONES
    \draw[->, thick] (pipeline) -- (woe);
    \draw[->, thick] (woe) -- (ca);
    \draw[->, thick] (ca) -- (ga);
    \draw[->, thick] (utils) -- (pipeline);
    \draw[->, thick] (ca) -- (eval);
    \draw[->, thick] (ga) -- (eval);
  \end{tikzpicture}
  \end{center}
  
  \vspace{0.3cm}
  \begin{columns}
    \begin{column}{0.5\textwidth}
      \textbf{Métricas de Código:}
      \begin{itemize}
        \item \textbf{2,847 líneas} Python documentadas
        \item \textbf{67 funciones} con tests unitarios
        \item \textbf{89\% cobertura} de testing
        \item \textbf{12 módulos} principales
      \end{itemize}
    \end{column}
    \begin{column}{0.5\textwidth}
      \textbf{Características:}
      \begin{itemize}
        \item Arquitectura modular extensible
        \item Configuración vía archivos YAML
        \item Documentación completa
        \item Framework reproducible
      \end{itemize}
    \end{column}
  \end{columns}
\end{frame}

\section{Estado del Arte Revisado}

\begin{frame}{Bibliografía Consolidada (\TotalRefs~fuentes)}
  \vspace{-0.4em}
  \begin{center}
    \begin{tikzpicture}[scale=0.8]
      \pie[
        sum=auto,
        after number=~refs,
        radius=2.5,
        text=legend,
  color={UNAMBlue!70,ForestGreen!70,Orange!70,UNAMRed!70,gray!50,gray!30},
        /tikz/every text node part/.style={align=left}
      ]{
        \RefsRS/Sensado remoto y DL,
        \RefsCA/AC y SLEUTH,
        \RefsCalib/Calibración (GA-PSO-MOO),
        \RefsWoE/Weight of Evidence,
        \RefsImpacto/Planeación e impacto,
        \RefsTools/Herramientas open-source
      }
    \end{tikzpicture}
  \end{center}
  
  \vspace{-0.3cm}
  \textbf{Hallazgos Clave:}
  \begin{itemize}
    \item \textbf{Brecha identificada:} Pocos estudios integran WoE + AC + AG de forma sistemática
    \item \textbf{Oportunidad:} Framework reproducible para ciudades intermedias mexicanas  
    \item \textbf{Tendencia:} Ecosistema open-source facilita transferencia tecnológica
  \end{itemize}
\end{frame}

\section{Validación y Resultados}

\begin{frame}{Protocolo de Validación Temporal Implementado}
  \begin{columns}
    \begin{column}{0.6\textwidth}
      \textbf{Estrategia de Validación:}
      \begin{enumerate}
        \item \textbf{Entrenamiento:} Datos 1984-2000 (16 años)
        \item \textbf{Calibración AG:} Optimización parámetros WoE-AC  
        \item \textbf{Validación:} Predicción 2000-2020 vs. observado
        \item \textbf{Métricas:} IoU, Kappa, Figure of Merit espacial
      \end{enumerate}
      
      \vspace{0.3cm}
      \textbf{Métricas Espaciales:}
      \begin{itemize}
        \item \textbf{IoU:} Intersección sobre Unión
        \item \textbf{Kappa:} Concordancia estadística  
        \item \textbf{FoM:} Figure of Merit para cambio urbano
        \item \textbf{Análisis residual:} Patrones error espacial
      \end{itemize}
    \end{column}
    \begin{column}{0.4\textwidth}
      \begin{center}
      \begin{tikzpicture}[scale=0.7]
        % Línea temporal
        \draw[thick] (0,0) -- (4,0);
        
        % Períodos
        \draw[ForestGreen, thick] (0,0) -- (2,0);
  \draw[UNAMRed, thick] (2,0) -- (4,0);
        
        % Etiquetas
        \node[below] at (1,0) {\scriptsize Entrenamiento};
        \node[below] at (3,0) {\scriptsize Validación};
        
        % Años
        \node[above] at (0,0) {\tiny 1984};
        \node[above] at (2,0) {\tiny 2000};
        \node[above] at (4,0) {\tiny 2020};
        
        % Proceso
        \node[below, font=\scriptsize] at (2,-0.8) {
          \begin{minipage}{3.5cm}
          \centering
          Calibrar parámetros AG\\
          ↓\\
          Predecir 2000-2020\\
          ↓\\
          Evaluar IoU, Kappa, FoM
          \end{minipage}
        };
      \end{tikzpicture}
      \end{center}
      
      \vspace{0.3cm}
      \begin{alertblock}{Precisión Esperada}
      \textbf{IoU >0.85} • \textbf{Kappa >0.80} • 
      \textbf{FoM >0.75} basado en literatura
      \end{alertblock}
    \end{column}
  \end{columns}
\end{frame}

\section{Contribuciones y Originalidad}

\begin{frame}{Contribuciones Metodológicas Innovadoras}
  \begin{enumerate}
    \item \textbf{Integración sistemática WoE-AC-AG:}
    \begin{itemize}
      \item Primera aplicación documentada de WoE en AC urbanos con AG
      \item Framework matemático riguroso para la integración
      \item Protocolo de calibración automática reproducible
    \end{itemize}
    
    \item \textbf{Validación temporal robusta (37 años):}
    \begin{itemize}
      \item Protocolo más extenso que literatura típica (2-5 años)
      \item Métricas espaciales especializadas para crecimiento urbano
      \item Análisis de degradación temporal sistemático
    \end{itemize}
    
    \item \textbf{Aplicación ciudades intermedias mexicanas:}
    \begin{itemize}
      \item Primer estudio AC detallado para Querétaro  
      \item Caracterización cuantitativa patrones crecimiento
      \item Base metodológica transferible otras ciudades similares
    \end{itemize}
  \end{enumerate}
\end{frame}

\section{Cronograma y Próximos Pasos}

\begin{frame}{Plan de Finalización - Cronograma Detallado}
  \begin{table}[h]
  \centering
  \footnotesize
  \begin{tabular}{p{2.5cm}p{3.5cm}p{4.5cm}p{1.5cm}}
  \toprule
  \textbf{Período} & \textbf{Actividad Principal} & \textbf{Entregables Específicos} & \textbf{Avance} \\
  \midrule
  \textcolor{ForestGreen}{\textbf{Dic 2025}} & Experimentos finales & 
  Calibración AG completa, mapas predicción vs. observación & \textcolor{ForestGreen}{\textbf{85\%}} \\
  \midrule
  Ene 2026 & Análisis resultados & 
  15 figuras técnicas, métricas espaciales, análisis sensibilidad & 60\% \\
  \midrule  
  Feb 2026 & Documentación final & 
  Tesis completa 140 páginas, presentación defensa & 40\% \\
  \midrule
  Mar 2026 & Defensa de tesis & 
  Examen de grado, publicaciones académicas & 20\% \\
  \bottomrule
  \end{tabular}
  \end{table}
  
\end{frame}

\begin{frame}
  \vspace{0.3cm}
  \begin{columns}
    \begin{column}{0.5\textwidth}
      \textbf{Actividades Críticas Diciembre:}
      \begin{itemize}
        \item Completar calibración AG con datos Querétaro
        \item Generar mapas predicción 2000-2020
        \item Análisis métricas espaciales detallado
      \end{itemize}
    \end{column}
    \begin{column}{0.5\textwidth}
      \begin{alertblock}{Meta Defensa}
      \textbf{Marzo 2026} • Tesis 140 páginas • 
      Framework reproducible • Publicación Q1
      \end{alertblock}
    \end{column}
  \end{columns}
\end{frame}

\section{Impacto Esperado}

\begin{frame}{Impacto Multi-Dimensional del Proyecto}
  \begin{columns}[T]
    \begin{column}{0.33\textwidth}
      \textcolor{UNAMBlue}{\textbf{Científico-Académico}}
      \begin{itemize}
        \item Cierra brecha metodológica en calibración AC
        \item Framework WoE + AG + métricas espaciales  
        \item \textbf{2-3 publicaciones} Q1 proyectadas
        \item Código abierto reproducible
        \item Base para investigación futura
      \end{itemize}
    \end{column}
    \begin{column}{0.33\textwidth}
      \textcolor{ForestGreen}{\textbf{Social-Aplicado}}
      \begin{itemize}
        \item Apoyo planeación urbana municipal
        \item Herramienta políticas desarrollo sostenible
        \item Evaluación impacto ambiental
        \item Transparencia toma decisiones  
        \item Transferible otras ciudades mexicanas
      \end{itemize}
    \end{column}
    \begin{column}{0.33\textwidth}
      \textcolor{Orange}{\textbf{Tecnológico}}
      \begin{itemize}
        \item Framework DEAP + Python escalable
        \item Integrable SIG modernos
        \item Procesamiento GPU/cluster
        \item Interfaz usuario planificadores
        \item Open-source ecosystem
      \end{itemize}
    \end{column}
  \end{columns}
\end{frame}


\section{Conclusiones}

\begin{frame}{Síntesis de Logros y Perspectivas}
  \begin{columns}
    \begin{column}{0.6\textwidth}
      \textbf{Logros Principales Noviembre 2025:}
      \begin{itemize}
        \item \textcolor{ForestGreen}{\textbf{$\checkmark$}} Dataset completo 37 años procesado (420 MB)
        \item \textcolor{ForestGreen}{\textbf{$\checkmark$}} Pipeline clasificación >99\% precisión  
        \item \textcolor{ForestGreen}{\textbf{$\checkmark$}} Sistema WoE completamente funcional
        \item \textcolor{UNAMBlue}{\textbf{$\checkmark$}} Framework AC-AG 95\% implementado
        \item \textcolor{ForestGreen}{\textbf{$\checkmark$}} Tesis 85\% redactada (140 páginas)
        \item \textcolor{ForestGreen}{\textbf{$\checkmark$}} Arquitectura software robusta (2,847 líneas)
      \end{itemize}
    \end{column}
    \begin{column}{0.4\textwidth}
      \textbf{Perspectiva Finalización:}
      \begin{itemize}
        \item Cronograma viable Marzo 2026
        \item Contribución metodológica sólida  
        \item Impacto científico y social
        \item Base para investigación futura
      \end{itemize}
      
      \vspace{0.3cm}
      \begin{alertblock}{Estado General}
      \textcolor{ForestGreen}{\textbf{85\% Completado}} \\
      En fase final experimentación
      \end{alertblock}
    \end{column}
  \end{columns}
  
  \vspace{0.4cm}
  \begin{center}
  \Large
  \textcolor{UNAMBlue}{\textbf{Framework WoE-AC-AG: Metodología sólida, implementación robusta, 
  resultados prometedores para modelado urbano en México}}
  \end{center}
\end{frame}

% REFERENCIAS
\appendix
\begin{frame}{Referencias Clave Expandidas}
  \footnotesize
  \textbf{Fundamentos Metodológicos:}
  \begin{itemize}
    \item Clarke \& Gaydos (1997); White \& Engelen (1997); Wolfram (1984) - \emph{Fundamentos AC}
    \item Almeida et al. (2003); Li \& Dang (2013); Wang et al. (2021) - \emph{AC + Optimización}
    \item Bonham-Carter (1994); Agterberg \& Cheng (2002) - \emph{Weight of Evidence}
  \end{itemize}
  
  \textbf{Datos y Procesamiento:}  
  \begin{itemize}
    \item Herold et al. (2003); Gong et al. (2013); Ma et al. (2019) - \emph{Sensado remoto urbano}
    \item Kadhim et al. (2018); Zhang et al. (2020) - \emph{Deep learning clasificación}
  \end{itemize}
  
  \textbf{Herramientas y Frameworks:}
  \begin{itemize}
    \item DEAP (Fortin et al., 2012); Repast (North et al., 2013) - \emph{Optimización evolutiva}  
    \item UrbanSim (Waddell, 2002); SLEUTH (Clarke, 2007) - \emph{Modelado urbano}
  \end{itemize}
  
  \textbf{Contexto Regional:}
  \begin{itemize}
    \item Angel et al. (2016); Seto et al. (2011, 2012); UN-Habitat (2020) - \emph{Crecimiento urbano global}
  \end{itemize}
\end{frame}

\begin{frame}[plain]
  \Huge ¡Gracias!
  
  \vspace{0.5cm}
  \Large Preguntas y Comentarios
  
  \vspace{0.3cm}
  \normalsize
  \textcolor{gray}{
  Rodrigo Moreno López \\
  Maestría en Ciencias de la Computación \\
  Universidad Nacional Autónoma de México
  }
\end{frame}

\end{document}